\documentclass[10pt,a4paper]{amsart}
\usepackage[margin=19mm]{geometry}
\usepackage{amsmath,amssymb,mathtools}
\usepackage[scr=rsfs,cal=euler]{mathalfa}
\usepackage{xfrac}
\usepackage[unicode,hidelinks,bookmarks=false]{hyperref}
\newcommand{\ie}{i.e.\ }
\newcommand{\cf}{cf.\ }
\newcommand{\resp}{respectively\ }
\newcommand{\Subsection}{\subsection}
\providecommand{\nobreakdash}{\nobreak\hbox{-}}
\setlength{\emergencystretch}{2em}
\multlinegap0pt
\usepackage{amssymb}
\usepackage{mathtools}
\usepackage{xcolor}
\newtheorem{res}{Result}
\newtheorem{thm}{Theorem}
\newcommand{\FF}{\mathbb F}
\newcommand{\ms}{\mathcal S}
 \DeclareMathOperator{\pr}{pr}
\newcommand{\sett}[2]{ \left\{ #1 \, \, || \, \, #2 \right \} }
\title{Points below a parabola in affine planes of prime order}
\author{Sam Adriaensen, Zsuzsa Weiner}
\date{Source: arXiv:2411.19202v2 (CC BY 4.0)}
\begin{document}
\maketitle

\noindent\emph{Source and licence.} Points below a parabola in affine planes of prime order. Version arXiv:2411.19202v2 (CC BY 4.0), distributed by arXiv under the Creative Commons Attribution 4.0 International licence, http://creativecommons.org/licenses/by/4.0/, as recorded in the arXiv licence metadata for this exact version and shown on the abstract page https://arxiv.org/abs/2411.19202v2. Full licence notice: \texttt{licenses/arxiv-2411.19202-CC-BY-4.0.txt}.\par\smallskip

\noindent\textit{Editorial scope.} Counted unit: the article's thm Thm:Parabola (source0.tex lines 298--311). The article's complete original proof of it is delivered with it (source0.tex lines 313--382, one \texttt{proof} environment of 70 lines). No mathematics is written, repaired or re-derived: the statement and the proof are the article's own lines, in the article's own notation, and no retained line is substituted or re-flowed.  Retained uncounted as the source-local context the proof needs: lines 230--236; lines 240--246.  Nothing was omitted from any retained span, so the package carries no omission record.  No retained line contains a citation, so the wrapper carries no bibliography and no bibliography entry.  No retained line contains a figure environment, an image-inclusion command or drawing code, so no image is delivered and the package declares no asset dependency.  Editorial formatting only, in addition: an AMS \texttt{amsart} preamble that restates the article's own declarations and macros used by the retained lines, namely the environments \texttt{res}, \texttt{thm} on separate counters, together with the standard packages those lines need.  The retained environments print in this document as Res 1, Theorem 1 in order of appearance, whereas the article prints the same items with the numbering of its own full document.  The complete native source of the article as distributed in the arXiv e-print for this exact version is bundled unchanged as \texttt{sources/169293/source0.tex}.
\begin{res}[P\'olya-Vinogradov inequality]
 \label{Res:PolyaVinogradov}
 For any two integers $a, b$,
 \[
  \left| \sum_{x=a}^b \chi(x) \right| \leq \sqrt p \ln(p).
 \]
\end{res}

\begin{res}
 \label{Res:Csikvari}
 There exists an element $b \in \FF_p$ such that
 \[
  \left| \sum_{x=1}^b \chi(x) \right| \geq \frac1{2\pi} \sqrt p.
 \]
\end{res}

\begin{thm}
 \label{Thm:Parabola}
 Take a quadratic polynomial $f(X) = \alpha X^2 + \beta X + \gamma$, $\alpha \neq 0$ over $\FF_p$, with $p > 2$ prime.
 Consider the set
 \[
  \ms = \sett{(x,y) \in \FF_p^2}{ y < f(x)}
 \]
 of points below the parabola $Y = f(X)$.
 For any $d \in \FF_p^*$ and $b \in \FF_p$,
 \[
  \pr_{\ms,d}(b) = \pr_{\ms,1}\left( b - \left(\beta-\frac{d+1}2\right)\frac{d-1}{2\alpha} \right).
 \]
 Moreover, the image of $\pr_{\ms,d}$ is an interval whose length lies between $\frac{\sqrt p}{2 \pi}$ and $\sqrt p \ln(p)$.
\end{thm}

\begin{proof}
 We need to prove that for any value of $b$,
 \[
  \pr_{\ms,d}\left( b + \left(\beta-\frac{d+1}2\right)\frac{d-1}{2\alpha} \right)
 \]
 is independent of the choice of $d \in \FF_p^*$.
 This integer equals the number of solutions to the inequality
 \[
  dX + b + \left(\beta-\frac{d+1}2\right)\frac{d-1}{2\alpha} < \alpha X^2 + \beta X + \gamma.
 \]
 We apply the linear transformation $X' = X - \frac{d-1}{2\alpha}$ to see that we may equivalently count the number of solutions to the inequality
 \begin{equation}
  \label{Eq:X'}
  X' + b + \delta_d(X') < \alpha X'^2 + \beta X' + \gamma + \delta_d(X'),
 \end{equation}
 with $\delta_d(X') = (d-1) X' + \frac{(d-1)^2}{4\alpha} + \beta \frac{d-1}{2\alpha}$.
 
 Let $N_{d,b}$ denote the set of solutions to Equation (\ref{Eq:X'}).
 Then we need to prove that $|N_{d,b}|$ is independent of the choice of $d \in \FF_p^*$.
 Since we know that $\sum_{b \in \FF_p} |N_{d,b}| = |\ms|$ for any $d$, it actually suffices to prove that for any $b$,
 \[
  |N_{d,b+1}|-|N_{d,b}| = |N_{d,b+1} \setminus N_{d,b}| - |N_{d,b} \setminus N_{d,b+1}|
 \]
 is independent of $d \in \FF_p^*$.
 Note that
 \begin{align*}
  x \in N_{d,b+1} \setminus N_{d,b} && \iff && x+b+\delta_d(x)=-1 && \text{ and } && f(x) + \delta_d(x) \neq 0, \\
  x \in N_{d,b} \setminus N_{d,b+1} && \iff && x+b+1=f(x) && \text{ and } && f(x) + \delta_d(x) \neq 0.
 \end{align*}
 Let $k$ denote the number of solutions to $X+b+1 = f(X)$.
 Since $d \neq 0$, there is a unique solution $x_0$ to $X + b + \delta_d(X) = -1$.
 
 First suppose that $x_0$ is a solution to $f(X) + \delta_d(X) = 0$.
 Then it is also a solution to $X+b+1 = f(X)$, hence
 \begin{align*}
  |N_{d,b+1} \setminus N_{d,b}| = 0, &&
  |N_{d,b} \setminus N_{d,b+1}| = k-1.
 \end{align*}

 If $x_0$ is not a solution to $f(X) + \delta_d(X) = 0$, then
 \begin{align*}
  |N_{d,b+1} \setminus N_{d,b}| = 1, &&
  |N_{d,b} \setminus N_{d,b+1}| = k.
 \end{align*}
 To see that the latter equality holds, we need to prove that no solution $x$ of $X+b+1 = f(X)$ satisfies $f(X) + \delta_b(X) = 0$.
 If that were the case, then $x + b + \delta_d(x) = -1$, which implies that $x = x_0$ and contradicts that $f(x_0) + \delta_d(x_0) \neq 0$.

 We conclude that in both cases
 \[
  |N_{d,b+1} \setminus N_{d,b}| - |N_{d,b} \setminus N_{d,b+1}| = 1-k,
 \]
 which is indeed independent of the choice of $d \in \FF_p^*$.

 Note that $k$ is the number of roots of the quadratic equation $f(X) - X - b - 1 = 0$.
 Therefore, $k$ depends on the quadratic character of the discriminant $(\beta-1)^2 + 4\alpha(b+1-\gamma)$.
 More specifically,
 \[
  1-k = -\chi( (\beta-1)^2 + 4\alpha(b+1-\gamma) ).
 \]
 This means $N_{d,b}$ and $N_{d,b+1}$ differ by $1-k$ and $|1-k| \leq 1$, hence that the image of $\pr_{\ms,d}$ consists of all integers in an interval.
 Moreover, for any $a,b \in \FF_p$,
 \begin{align*}
  \Big | |N_{d,b}| - |N_{d,a}| \Big |
  &= \left | \sum_{x = a+1}^b \chi( (\beta-1)^2 + 4\alpha(x-\gamma) ) \right| \\
  &= |\chi(4 \alpha)| \left | \sum_{x = a+1}^b \chi \left( \frac{(\beta-1)^2}{4\alpha} + x-\gamma \right) \right|
  = \left | \sum_{x = \frac{(\beta-1)^2}{4\alpha} + a+1-\gamma}^{\frac{(\beta-1)^2}{4\alpha} + b-\gamma} \chi \left( x \right) \right|
 \end{align*}
 Therefore, the length of the interval formed by the image of $\pr_{\ms,d}$ equals the maximum value of the above sum.
 We know from Results \ref{Res:PolyaVinogradov} and \ref{Res:Csikvari} that this maximum value is between $\frac{\sqrt p}{2 \pi}$ and $\sqrt p \ln(p)$.
\end{proof}

\end{document}
